\section*{Introdução}

O pêndulo simples é um sistema físico constituído por uma massa puntiforme suspensa por um fio inextensível e de massa desprezível, capaz de oscilar sob a ação da gravidade. Para pequenas amplitudes angulares, o movimento do pêndulo pode ser aproximado por um movimento harmônico simples, permitindo relacionar o período de oscilação com o comprimento do fio e a aceleração da gravidade, conforme a equação ~\eqref{eq:periodo_pendulo}.

\begin{equation}
T = 2\pi \sqrt{\frac{L}{g}}
\label{eq:periodo_pendulo}
\end{equation}

Elevando ambos os lados da equação ~\eqref{eq:periodo_pendulo} ao quadrado, obtemos ~\eqref{eq:periodo_quadrado}.

\begin{equation}
T^2 = \left(2\pi \sqrt{\frac{L}{g}}\right)^2
\label{eq:periodo_quadrado}
\end{equation}
